\chapter{Prólogo}
\markboth{Prólogo}{Prólogo}

\begin{documento}[Compendio Universal de Ciencia Archivística · 4.ª edición, Ceres, 2231]
\textsc{Calderón, Inés} (Lima, Tierra, 2131 -- Ceres, 2209). Matemática de sistemas. Se la recuerda, sobre todo, por haber resuelto la llamada \emph{Crisis del Cartel} de 2187, que había paralizado durante once días el Archivo Central de Ceres, y por la solución que lleva su nombre en los manuales elementales, aunque ella siempre se negó a llamarla así.

Antes de 2187, el Archivo cumplía su Reglamento a costa de una de sus tres Reglas. O se cerraban los carteles con llave, y los lectores esperaban; o se dejaban sin fundir los cristales viejos durante semanas, y la Bóveda se llenaba de basura; o se aceptaba, en silencio, el riesgo de que un lector llegara a un nicho cuyo cristal ya no existía. La solución de Calderón permitió, por primera vez, cumplir las tres a la vez.

Es característico del personaje que, preguntada al final de su vida por su mayor aportación a la ciencia, respondiera: «Ninguna. Solo miré el cartel, cuando todos miraban la hoja». Y es característico de la historia de la ciencia que, años después, se descubriera que no había sido la primera en mirarlo (véase el \emph{Epílogo} de esta entrada).

\textsc{Véase también:} \emph{Primera Regla}; \emph{ticket prepagado}; \emph{zona roja}; \emph{salida temprana}; \emph{Siete (robot)}.
\end{documento}

\bigskip

\inicio{E}{ste} libro cuenta cómo ocurrió aquello. Lo cuenta, además, con el máximo de verdad que permiten las novelas: las personas y los robots son inventados, Ceres no tiene ningún Archivo Central y ningún carguero llamado \emph{Perséfone} estuvo nunca a punto de estrellarse por culpa de un inventario de trigo. Pero el problema es real, la solución es real y cada número que aparece en la pizarra de Inés Calderón se puede comprobar con lápiz y papel.

El problema existe en casi todos los ordenadores que usamos. Cada vez que un programa con varios hilos comparte un dato y lo cuenta para saber cuándo tirarlo, se arriesga a lo que en estas páginas llamaremos \emph{la carrera del cartel}: un hilo mira dónde está el dato; otro, en ese mismo instante, lo tira; el primero llega y encuentra en su lugar otra cosa. Los informáticos lo llaman \emph{carrera read-reclaim} y su consecuencia, \emph{use-after-free}. Es una de las fuentes de fallos de seguridad más frecuentes que existen.

En 2024, cuatro investigadores ---Ivo Gabe de Wolff, Daniel Anderson, Gabriele K.\ Keller y Aleksei Seletskiy--- publicaron una solución rápida, sencilla y \emph{sin esperas}. Lo que sigue es esa solución, contada como un misterio. El lector que solo quiera la historia puede saltarse las pizarras. El que quiera entenderla a fondo encontrará, al final, un diccionario que traduce cada objeto de Ceres a su nombre técnico, los algoritmos completos y la lista de lo que es invención y lo que no.

Empecemos, pues, como empiezan estas cosas: con una nave que estaba a punto de chocar.
